
\documentclass[12pt,a4paper]{article}
\usepackage[utf8]{inputenc}
\usepackage[spanish]{babel}
\usepackage{geometry}
\geometry{margin=2.5cm}
\usepackage{amsmath,amssymb}
\usepackage{graphicx}
\usepackage{hyperref}

\title{\textbf{Informe Académico Automatizado: Analítica: ¿Cuántas personas llegaron al restaurante, cu...}}
\author{Sistema Local RAG - IDEAD}
\date{\today}

\begin{document}

\maketitle

\begin{abstract}
Reporte enfocado en: ¿Cuántas personas llegaron al restaurante, cuál fue su distribución espacial y qué zonas registran mayor afluencia?
\end{abstract}

\section{Introducción y Contexto}
Este documento ha sido generado automáticamente como parte de la recopilación bibliográfica sobre \textbf{Minería de Datos Aplicada a Flujos de Video y Ocupación}, integrando motores de búsqueda académica y bases de datos vectoriales locales.

\section{Desarrollo Teórico}
Planteamiento del problema: ¿Cuántas personas llegaron al restaurante, cuál fue su distribución espacial y qué zonas registran mayor afluencia?\newline\newlineNotas operativas del analista: Análisis de flujos y ocupación espacial basado en minería de datos local mediante CRISP-DM.\newline\newlineResultados cuantitativos y analíticos:\newline Metodología CRISP-DM aplicada sobre un total de 500 registros. Distribución detectada por zonas: Zona Mesas Principal: 259 registros, Zona Ingreso/Barra: 241 registros.\newline\newlineConclusiones: El procesamiento de flujos de video mediante bases de datos analíticas (DuckDB) permite dar respuesta directa a los objetivos planteados, optimizando la gestión del establecimiento.

\section{Referencias}
\begin{enumerate}
    \item Recopilado mediante cadenas booleanas optimizadas para investigación en ingeniería.
\end{enumerate}

\end{document}
